\item De acuerdo con la física clásica, una carga $e$ móvil con una aceleración $a$ emite energía a una tasa de:
$$ \frac{dE}{dt} = -\frac{2}{3} \frac{e^2 a^2}{4\pi\varepsilon_0 c^3} $$
\begin{enumerate}
    \item Demuestre que un electrón en un átomo de hidrógeno clásico se mueve en espiral hacia el núcleo con una rapidez radial dada por:
    $$ \frac{dr}{dt} = -\frac{2}{3} \left( \frac{e^4}{m_e^2 c^3} \right) \left( \frac{1}{4\pi\varepsilon_0} \right)^2 \frac{1}{r^2} $$
    \item Determine el intervalo de tiempo al final del cual el electrón alcanzará $r = 0$, empezando desde $r_0 = 2.00 \times 10^{-10}\text{ m}$.
\end{enumerate}